\subsection{Normal-crossings fibres with arbitrary multiplicities}

In this subsection a normal-crossings surface is allowed to have self-intersections: distinct local branches need not belong to distinct global irreducible components. We count a triple point once in the surface, even when its three branches lie on the same global component. Euler characteristics of noncompact strata mean compactly supported Euler characteristics; all compact spaces have their usual Euler characteristic. All cohomology classes used numerically have rational coefficients.

\begin{statement}{C.2.NC.1}{Known, with derivation}{The local-to-global Euler formula with multiplicities and self-intersections}
Let $g:U\to\Delta$ be a proper holomorphic map from a smooth complex threefold to a disc, smooth over $\Delta^*$, whose central fibre is $\sum_i m_iS_i$ and whose reduced support $S$ has normal crossings. Let $\nu:\coprod_i T_i\to S$ be normalization and let $D_i\subset T_i$ be its reduced conductor divisor. For a sufficiently small nonzero $t$,
\[
 e(U_t)=\sum_i m_i e(T_i\setminus D_i).
 \tag{C.NC.1}
\]
The formula does not require a Kähler total space, a polarization, uniform multiplicities, or simple normal crossings.
\end{statement}
\begin{proof}
We spell out the normal-crossings form of the construction in \cite[\S2.1--2.2, pp.~237--241]{ACampo1975}. Replace each local complex normal coordinate by its modulus and argument. In a chart with $r$ local branches, absorb a local root of the unit in the equation of $g$ into one coordinate, so that
\[
 g=z_1^{a_1}\cdots z_r^{a_r},\qquad
 z_j=\rho_j e^{\mathrm i\theta_j},\qquad 1\le r\le3.
\]
The resulting real oriented blowup $Z\to U$ has a boundary with corners over $S$. This construction glues even when branches are permuted: a change of defining equation multiplies a normal coordinate by a nowhere-zero holomorphic function, which multiplies its modulus by a positive function and translates its argument; permutation of the branches permutes the corresponding factors. Thus the construction uses the local branches themselves and does not require a global ordering or a global separation of the branches.

The lift of $g$ to the real oriented blowup of the disc is given in these coordinates by
\[
 (\rho_j,\theta_j)\longmapsto
 \left(\prod_j\rho_j^{a_j},\ \sum_j a_j\theta_j\pmod {2\pi}\right).
 \tag{C.NC.2}
\]
This map is locally a topological projection, including along the corners. For the radial variables this can be checked explicitly: put $u_j=\rho_j-\rho_r$ for $j<r$ and $v=\prod_j\rho_j^{a_j}$. Given $(u_1,\ldots,u_{r-1},v)$, the missing variable $\lambda=\rho_r$ is the unique solution, for $v>0$, of
\[
 \lambda^{a_r}\prod_{j<r}(u_j+\lambda)^{a_j}=v,
 \quad \lambda\ge\max(0,-u_1,\ldots,-u_{r-1});
\]
for $v=0$ it is the endpoint of that interval. Strict monotonicity for positive $v$ and continuity at the endpoint give the inverse radial chart. The angular coordinate in (C.NC.2) is a submersion in local real angular charts. Properness and the resulting local product charts give the boundary fibration used in the cited construction; its fibre over an argument $\theta$ is homeomorphic to $U_t$ for $\arg t=\theta$. This is exactly the comparison proved in the lemma of \cite[\S2.2]{ACampo1975}; the preceding chart calculation also shows that its proof is unaffected by branch permutations.

The projection of that boundary fibre to $S$ has, over a point with $r$ local branches, fibre
\[
 \left\{(u_1,\ldots,u_r)\in(S^1)^r:
             \prod_j u_j^{a_j}=e^{\mathrm i\theta}\right\}.
 \tag{C.NC.3}
\]
This is the disjoint union of $\gcd(a_1,\ldots,a_r)$ copies of $(S^1)^{r-1}$: an integral change of basis taking the row $(a_1,\ldots,a_r)$ to $(\gcd(a_j),0,\ldots,0)$ gives the asserted homeomorphism. It has Euler number zero for $r\ge2$ and $a_1$ for $r=1$.

Choose a finite triangulation of the compact analytic strata of $S$, refined so that the normal ray bundles and branch coverings trivialize on every open cell. Over each cell (C.NC.3) is a locally trivial bundle. Additivity over the cells and multiplicativity of compactly supported Euler characteristic give zero for every multiple-branch stratum. On the one-branch stratum of $S_i$ the fibre has $m_i$ points. Normalization identifies that stratum with $T_i\setminus D_i$, giving (C.NC.1). This also proves directly which local contributions enter the global formula.
\end{proof}

\begin{statement}{C.2.NC.2}{Proved}{The normalized double curves and the coherent Euler formula}
Let $S$ be a connected compact normal-crossings divisor in a smooth compact complex threefold $Y$ with $H^2(Y,\Q)=0$. Write $\nu:\coprod_iT_i\to S$ for normalization, $D_i\subset T_i$ for its conductor, $C_\alpha$ for the connected components of the normalization of the double locus, and $\tau$ for the number of triple points. Let $R_\alpha$ be the reduced divisor on $C_\alpha$ of preimages of triple points, counted separately for the three double-curve branches at each triple point. Then
\[
 e(C_\alpha)=\deg R_\alpha,\qquad
 \sum_\alpha e(C_\alpha)=3\tau,
 \qquad
 \sum_i\chi(T_i,\OO_{T_i})=\frac{\tau}{2}.
 \tag{C.NC.4}
\]
In particular every normalized double curve has genus at most one, and $\tau$ is even. The remaining topological identities are
\[
\begin{aligned}
 \sum_i e(D_i)&=2\sum_\alpha e(C_\alpha)-3\tau,\\
 e(S)&=\sum_i e(T_i)-\sum_\alpha e(C_\alpha)+\tau.
\end{aligned}
 \tag{C.NC.5}
\]
\end{statement}
\begin{proof}
Adjunction and the local residue comparison give a canonical holomorphic line-bundle isomorphism
\[
 \nu^*\omega_S|_{T_i}\simeq\omega_{T_i}(D_i),
 \qquad \omega_S=(\omega_Y\otimes\OO_Y(S))|_S.
 \tag{C.NC.6}
\]
For example, the residue of $dx\wedge dy\wedge dz/(xy)$ on the branch $x=0$ is $\pm dy\wedge dz/y$, and at $xyz=0$ it is $\pm dy\wedge dz/(yz)$. These expressions agree under changes of local equations by the usual determinant and unit factors of the displayed line bundles, proving the comparison. Since both line bundles on the right in the expression for $\omega_S$ come from $Y$, its first Chern class and that of $\omega_{T_i}(D_i)$ vanish rationally.

Normalize every irreducible curve in every $D_i$. The resulting disjoint union $\widehat D$ maps to $\coprod_\alpha C_\alpha$ by an étale map of total degree two over each $C_\alpha$. Away from a triple point this follows from the two branches of $xy=0$. At a triple point, each of the three normalized double branches has exactly two preimages, and their maps in the remaining coordinate are local isomorphisms. This proves étaleness at the missing points as well. The covering over a given $C_\alpha$ may be connected; no splitting is assumed.

Adjunction for a nodal divisor on a smooth surface, pulled back to $\widehat D$, identifies the restriction of $\omega_{T_i}(D_i)$ with the canonical line of that normalized conductor curve twisted by all its preimages of nodes of $D_i$. Such nodes are exactly the preimages of triple points of $S$. This is the elementary local dualizing-form identity $dw/w$ on each branch of a node. The resulting degree is zero by (C.NC.6). Since the covering of $C_\alpha$ is étale and its marked divisor is the pullback of $R_\alpha$, it follows that
$\deg\omega_{C_\alpha}+\deg R_\alpha=0$.
Thus $e(C_\alpha)=\deg R_\alpha$. There are exactly three marks over every triple point, even when two or three belong to one $C_\alpha$, proving their summed formula.

For the coherent calculation, the alternating branch-resolution sequence is
\[
 0\longrightarrow\OO_S\longrightarrow\nu_*\OO_{\coprod T_i}
 \longrightarrow\bigoplus_\alpha(j_\alpha)_*\varepsilon_\alpha
 \longrightarrow\bigoplus_{p\in S^{[3]}}\C_p
 \longrightarrow0.
 \tag{C.NC.7}
\]
Here $j_\alpha:C_\alpha\to S$ is the finite normalization map of that double curve, and $\varepsilon_\alpha$ is the sign line of its two-sheeted conductor covering. Locally order the branches and send $(f_i)$ to $(f_j-f_i)_{i<j}$; the next map is the alternating sum of the three values at a triple point. The sequence is exact at $xy=0$ because a pair of functions glues precisely when its restrictions agree, and at $xyz=0$ because pairwise differences satisfy exactly the single alternating relation. Equivalently, compare coefficients of convergent power series on the three coordinate planes; equality on their coordinate axes is exactly the condition for these series to come from one series modulo $xyz$. Changing a local branch ordering multiplies the appropriate difference by $-1$, producing $\varepsilon_\alpha$. At a triple point this sign convention only chooses a basis of the one-dimensional last term. Hence the sequence is globally well-defined with the indicated local systems, including self-intersections.

The sign line satisfies $\varepsilon_\alpha^{\otimes2}\simeq\OO_{C_\alpha}$, so it has degree zero and $\chi(C_\alpha,\varepsilon_\alpha)=1-g(C_\alpha)=e(C_\alpha)/2$. The divisor sequence on $Y$ and Hirzebruch--Riemann--Roch give $\chi(S,\OO_S)=0$, since $c_1(\OO_Y(S))=0$ rationally. Taking Euler characteristics in (C.NC.7) now gives
\[
 0=\sum_i\chi(T_i,\OO_{T_i})-\frac12\sum_\alpha e(C_\alpha)+\tau,
\]
which proves the last equality in (C.NC.4).

Finally $\widehat D$ has Euler number $2\sum_\alpha e(C_\alpha)$. Each triple point produces three nodes in $\coprod D_i$, and identifying the two preimages of each node subtracts one. This proves the first identity in (C.NC.5). The double locus in $S$ has normalization $\coprod C_\alpha$ and three preimages over each triple point, so its Euler number is $\sum_\alpha e(C_\alpha)-2\tau$. Normalization is an isomorphism off the conductor. Additivity of Euler characteristic therefore gives the second identity in (C.NC.5).
\end{proof}

\begin{statement}{C.2.NC.3}{Proved}{Uniform multiple normal-crossings fibres}
Let $g:Y\to C$ be a proper holomorphic map from a connected smooth compact complex threefold with $H^2(Y,\Q)=0$ to a smooth compact curve, with connected fibres and smooth general fibre of Euler characteristic zero. A fibre $mS$ with reduced normal-crossings support has
\[
 e(S)=\tau(S)\ge0.
 \tag{C.NC.8}
\]
The normalized base change $t=s^m$ is finite étale of degree $m$ over a neighbourhood of $S$ and has reduced normal-crossings central fibre $S'$. Its possibly disconnected components do not affect the comparisons
$e(S')=m e(S)$ and $\tau(S')=m\tau(S)$.
\end{statement}
\begin{proof}
The local formula C.2.NC.1 gives
\[
 0=e(Y_t)=m\sum_i e(T_i\setminus D_i).
\]
By (C.NC.5),
\[
 e(S)-\sum_i e(T_i\setminus D_i)
 =\sum_\alpha e(C_\alpha)-2\tau=\tau.
\]
This proves (C.NC.8) directly, with no base-change assumption.

To verify the additional comparison, locally write $t=u(z_1\cdots z_r)^m$ with $u$ nowhere zero. The normalization of $s^m=u(z_1\cdots z_r)^m$ is given by
$w^m=u$ and $s=wz_1\cdots z_r$.
The derivative of $w^m-u$ with respect to $w$ is nonzero because $w$ is a unit. Thus its map to the original chart is étale of degree $m$, and the lifted central fibre has equation $z_1\cdots z_r=0$. These local normalizations glue by the universal property of normalization; their degree and étaleness are local properties. Restriction is therefore an étale degree-$m$ covering $S'\to S$, and each triple point has exactly $m$ preimages. A lifted triangulation proves the Euler-number comparison. If the cover has several connected components, the degrees and Euler numbers add over them. The two displayed comparisons remain valid without any connectedness assertion about that cover.
\end{proof}

\begin{statement}{C.2.NC.4}{Proved}{Reducible normal-crossings fibres with unequal multiplicities}
Under the hypotheses on $g:Y\to C$ in C.2.NC.3, let a fibre be $\sum_i m_iS_i$ with connected reducible normal-crossings support $S$. Its multiplicities may be unequal. Then
\[
 \tau(S)=0,\qquad e(S)=0.
 \tag{C.NC.9}
\]
Every normalized conductor is a disjoint union of smooth elliptic curves, and every normalization component has Euler number zero. More generally, any singular normal-crossings fibre support under these hypotheses has either no triple point or exactly two triple points; in the latter case it is irreducible and its normalization is rational.
\end{statement}
\begin{proof}
First consider a normalization component $T_i$ whose conductor $D_i$ is nonempty. If $T_i$ is Kähler and has nonnegative Kodaira dimension, choose a nonzero section of $\omega_{T_i}^{\otimes n}$ and denote its effective zero divisor by $E$. The effective divisor $E+nD_i$ has rational Chern class zero by (C.NC.6). Its integral against a Kähler form is positive because $D_i$ is nonempty, whereas its cohomology class gives integral zero. This is impossible. A Kähler $T_i$ is consequently rational or a blown-up ruled surface. For a ruling $T_i\to B$ and a general ruling fibre $F$, one has $K_{T_i}\cdot F=-2$, hence $D_i\cdot F=2$. Some conductor curve therefore dominates $B$. Its normalization has genus at most one by C.2.NC.2, so Riemann--Hurwitz forces $g(B)\le1$. The Kähler possibilities have
$\chi(\OO_{T_i})=1$ in the rational case and $0$ in the elliptic ruled case; their Euler numbers are nonnegative.

For a non-Kähler $T_i$, the Enriques--Kodaira table \cite[Theorem~2.2]{Peters2021} leaves class VII, Kodaira and properly elliptic minimal models. Class VII has $c_1^2=-b_2$ and $c_2=b_2$, so Noether's formula gives $\chi(\OO)=0$. Kodaira surfaces have both Chern numbers zero. A minimal non-Kähler elliptic surface has topological Euler number zero by Kodaira's theorem in the precise form \cite[Theorem~3.2]{BMS2012}; its canonical square is zero by the same classification table, so Noether's formula again gives $\chi(\OO)=0$. Blowing up preserves the holomorphic Euler characteristic and adds one to the topological Euler characteristic. Thus every non-Kähler $T_i$ has $\chi(\OO_{T_i})=0$ and $e(T_i)\ge0$. This derives all numerical classification facts used here from the cited statements.

A rational normalization component cannot meet a different global component of $S$. Indeed, for $j\ne i$, pull back the canonical section of $\OO_Y(S_j)$ to $T_i$. It is a nonzero section whose zeros are precisely the inverse image of $S_i\cap S_j$. Its line bundle has zero rational first Chern class. For a smooth rational surface $T_i$, $H^1(T_i,\OO)=0$ and $H^2(T_i,\Z)$ is torsion-free, by the corresponding facts for $\PP^2$ or a Hirzebruch surface and the blowup formula. The exponential sequence therefore makes that pulled-back line bundle trivial. A nonzero section of a trivial line bundle on a connected compact surface is a nonzero constant and has no zeros. Hence $S_i\cap S_j$ is empty. This proof places no restriction on intersections between two local branches belonging to the same $S_i$.

The intersection graph of the global components of a connected reducible fibre is connected. Each component has nonempty conductor, and the preceding paragraph excludes rational normalizations. Thus every $\chi(\OO_{T_i})$ is zero, and (C.NC.4) gives $\tau=0$. Each $C_\alpha$ is now elliptic and has no marked point. The conductor covering is étale, so the normalized conductor components are elliptic; there are no nodes of $D_i$, since these would be preimages of triple points. Consequently $e(D_i)=0$. The weighted Euler formula now reads
\[
 0=e(Y_t)=\sum_i m_i e(T_i).
\]
All summands are nonnegative and every $m_i$ is positive, so each $e(T_i)=0$. Formula (C.NC.5) gives $e(S)=0$.

For a singular irreducible support the conductor is nonempty and the same classification argument gives $\chi(\OO_{T_1})\in\{0,1\}$, with value one exactly in the rational case. Formula (C.NC.4) then gives $\tau=0$ or $2$, respectively. This proves the final assertion and also proves that any triple point lies on an irreducible support with rational normalization.
\end{proof}

\begin{statement}{C.2.NC.5}{Proved}{Global positivity and the all-twist normalization obstruction in the normal-crossings class}
Let $g:Y\to C$ have the hypotheses of C.2.NC.3, and assume every fibre reduction has normal crossings; smooth reductions are included. Then
\[
 c_3(Y)=e(Y)=\sum_{c\in C}\tau((Y_c)_{\rm red})=2r\ge0,
 \tag{C.NC.10}
\]
where $r$ is the number of fibre supports with rational normalization. Every such support is irreducible, nonnormal, and has exactly two triple points. Equality $c_3(Y)=0$ holds exactly when every fibre support is free of triple points. If also $b_1(Y)=0$, then $r\in\{0,1\}$ and $b_3(Y)=2-2r$.

In the positive case there is an irreducible nonnormal reduced support $S$ with rational normalization $\eta:T\to S$ such that, for every $L\in\Pic(Y)$,
\[
 H^2\!\left(T,K_T\otimes\eta^*((K_Y^{-1}\otimes L)|_S)\right)\simeq\C,
 \qquad
 H^0\!\left(S,\Omega_Y^1|_S\otimes(K_Y\otimes L^{-1})|_S\right)\simeq\C.
 \tag{C.NC.11}
\]
The scheme-theoretic fibre may be a multiple of $S$; this conclusion does not assert that its multiplicity is one. By C.3.5, all these conclusions apply to every fibration under the CDP hypotheses whose fibre reductions have normal crossings, for every smooth-fibre type in C.3.1.
\end{statement}
\begin{proof}
There are finitely many critical values: the critical locus is analytic, its image is analytic by properness, and it does not contain a general point of the curve. Away from these values $g$ is a differentiable fibre bundle whose fibre has Euler number zero. Euler characteristic is additive over this open base and the finite exceptional set, so $e(Y)=\sum_c e((Y_c)_{\rm red})$. Nilpotents leave the underlying space unchanged. An irreducible support has a single multiplicity, so C.2.NC.3 applies; a reducible support is covered by C.2.NC.4. A smooth support has Euler number zero directly from (C.NC.1). Thus every summand equals its triple-point count. C.2.NC.4 identifies the positive summands and their value two. Gauss--Bonnet gives $c_3(Y)=e(Y)$. Poincaré duality gives $e(Y)=2-2b_1(Y)+2b_2(Y)-b_3(Y)$, proving the assertion about $r$ and $b_3$. C.3.5 proves the required zero Euler number for every smooth fibre under the CDP hypotheses, so that application uses no assumption that the general fibre is a torus.

For a positive support, every line bundle on $Y$ pulls back trivially to its smooth rational normalization: its rational first Chern class is zero, and the injectivity and torsion-freeness used in C.2.NC.4 apply. In particular $\eta^*((K_Y^{-1}\otimes L)|_S)$ is trivial. Serre duality on the connected smooth compact surface $T$ identifies the first group in (C.NC.11) with $H^0(T,\OO_T)^\vee\simeq\C$.

For the ambient group set $A=(K_Y\otimes L^{-1})|_S$ and $N^*=\OO_Y(-S)|_S$. Both $\eta^*A$ and $\eta^*N^*$ are trivial, even when the full fibre is $mS$. The injection of torsion-free differentials into normalization differentials gives
$H^0(S,\widetilde\Omega_S^1\otimes A)\hookrightarrow H^0(T,\Omega_T^1)=0$.
The canonical normal-crossings kernel isomorphism A.1.2 therefore gives
\[
 H^0(S,\Omega_Y^1|_S\otimes A)
 =H^0(T,\eta^*(N^*\otimes A))=\C.
\]
This proves the second equality. It establishes the universal normalization target requested in B.3 for the full normal-crossings class, including unequal multiplicities. It leaves the unrestricted nonordinary-fibre assertion to be decided separately.
\end{proof}
